% © 2026 Ing. David Jaroš — CC BY-NC-ND 4.0
%
% This work is licensed under a Creative Commons Attribution-NonCommercial-NoDerivatives
% 4.0 International License (CC BY-NC-ND 4.0).
%
% License History: Earlier drafts (up to v0.3) were released under CC BY 4.0.
% From v0.4 onward, all material is released under CC BY-NC-ND 4.0 to protect
% the integrity of the theoretical work during ongoing academic development.

\documentclass[12pt]{article}
\usepackage{amsmath,amssymb,amsthm}
\usepackage{geometry}
\usepackage{hyperref}
\usepackage{tcolorbox}
\usepackage{xcolor}
\geometry{margin=1in}

% Theorem-like environments
\newtheorem{definition}{Definition}
\newtheorem{proposition}{Proposition}
\newtheorem{theorem}{Theorem}
\newtheorem{remark}{Remark}
\newtheorem{conjecture}{Conjecture}

\title{ST-4: Connection Between Imaginary-Sector Invisibility\\
and CTC Causal Structure in UBT\\[0.5em]
\large\textsc{Speculative Extension — Invisibility Track}}
\author{David Jaroš}
\date{2026-03-03}

\begin{document}
\maketitle

\begin{tcolorbox}[colback=yellow!10, title=Status]
\textbf{SPECULATIVE EXTENSION} — This content is not part of
the canonical UBT framework.  It is an exploratory investigation
and has not been validated against the canonical field equations.
\end{tcolorbox}

\vspace{1em}
\noindent\textit{Output of Sub-task ST-4 from the Invisibility / Imaginary Sector
investigation track.  Cross-references:
\texttt{speculative\_extensions/causality/ctc\_conditions.tex} and
\texttt{speculative\_extensions/causality/complex\_time\_causal\_structure.tex}.
See also:
\texttt{st1\_imaginary\_object\_definition.tex},
\texttt{st2\_sector\_interaction.tex},
\texttt{st3\_stability.tex}.}

%----------------------------------------------------------------------
\section*{Abstract}
%----------------------------------------------------------------------
We examine whether the imaginary-sector invisibility mechanism
identified in ST-1–ST-3 is related to the Closed Timelike Curve (CTC)
causal structure analysed in the CTC investigation track.
The central question: \emph{Is an object in the imaginary sector
also outside the causal structure defined for CTCs?}
We argue that invisibility and CTCs share a common ingredient —
exit from the real metric sector $g_{\mu\nu}$ — but represent
\emph{distinct phenomena}: invisibility requires $\Theta_I$-dominated
energy-momentum with no contribution to $g_{\mu\nu}$, whereas CTCs
require $g_{\mu\nu}$ itself to admit closed causal curves.
The two effects can nevertheless coexist in the same UBT solution.
All conclusions are \textbf{speculative}.

%----------------------------------------------------------------------
\section{Summary of the CTC Causal Structure (from CTC Track)}
%----------------------------------------------------------------------

The CTC investigation (\texttt{ctc\_conditions.tex},
\texttt{complex\_time\_causal\_structure.tex}) established the following:

\begin{enumerate}
  \item The causal structure of UBT is determined primarily by the
    \emph{real} metric $g_{\mu\nu} = \operatorname{Re}(\mathcal{G}_{\mu\nu})$
    under the preferred Interpretation B (imaginary time as internal
    phase coordinate).

  \item The imaginary metric sectors $\{h_{\mu\nu}, a_{\mu\nu},
    b_{\mu\nu}\}$ contribute an effective stress-energy
    $T_{\mu\nu}^{(h)} = \operatorname{Re}(\partial\mathcal{L}/
    \partial\mathcal{G}^{\mu\nu})$ that can violate the weak energy
    condition, providing the mechanism for CTC-compatible geometries
    (Conjecture~1 in \texttt{ctc\_conditions.tex}).

  \item CTCs arise when the real metric $g_{\mu\nu}$ admits a Killing
    vector $\xi^\mu$ with compact orbits satisfying
    $g_{\mu\nu}\xi^\mu\xi^\nu < 0$ (the UBT-CTC condition).

  \item No explicit UBT solution with CTCs has been constructed;
    the CTC analysis identified mechanisms, not explicit solutions.
\end{enumerate}

%----------------------------------------------------------------------
\section{Summary of Invisibility (from ST-1--ST-3)}
%----------------------------------------------------------------------

\begin{enumerate}
  \item An imaginary-sector object $\Theta_I$ has
    $\operatorname{Re}(\mathcal{G}_{\mu\nu}[\Theta_I]) = 0$ (or
    near-zero) — it does not directly contribute to $g_{\mu\nu}$.

  \item Its energy-momentum tensor $T_{\mu\nu}[\Theta_I]$ is real
    and non-zero, so it gravitates and can generate lensing.

  \item At leading order, $\Theta_I$ does not couple to the
    electromagnetic field, making it electromagnetically invisible.

  \item Imaginary-sector objects are perturbatively stable, with
    long-lived or permanent invisibility depending on the strength
    of $U(1)$ breaking.
\end{enumerate}

%----------------------------------------------------------------------
\section{Comparison: Common Ingredient and Differences}
%----------------------------------------------------------------------

\begin{center}
\renewcommand{\arraystretch}{1.5}
\begin{tabular}{|p{4cm}|p{4.5cm}|p{4.5cm}|}
\hline
\textbf{Property} & \textbf{Invisibility ($\Theta_I$)} & \textbf{CTCs}\\
\hline
Role of $\operatorname{Im}(\mathcal{G}_{\mu\nu})$ &
  $\Theta_I$ lives here; no $g_{\mu\nu}$ contribution &
  $h_{\mu\nu}$ provides negative energy to tilt light cones\\
\hline
Effect on $g_{\mu\nu}$ &
  None directly (invisible to real metric) &
  Essential: CTCs defined by $g_{\mu\nu}$ admitting closed causal curves\\
\hline
Energy condition &
  $T_{\mu\nu}[\Theta_I]$ real, positive energy (Proposition 1, ST-3) &
  Negative effective energy density from $h_{\mu\nu}$ required\\
\hline
Observable signature &
  Gravitational lensing without visible source &
  CTC onset: modified Killing horizon radii\\
\hline
Stability &
  Perturbatively stable ($Q_I$ conserved) &
  Chronology protection may forbid stable CTCs\\
\hline
GR limit $\psi\to 0$ &
  $\Theta_I$ decouples, dark matter signal vanishes &
  Standard causal structure recovered\\
\hline
\end{tabular}
\end{center}

The \emph{common ingredient} is the imaginary sector of the biquaternionic
metric: both invisibility and CTCs arise from $\operatorname{Im}(\mathcal{G}_{\mu\nu})
\neq 0$.  However, they use this ingredient in opposite ways:
\begin{itemize}
  \item \textbf{Invisibility}: The imaginary sector \emph{contains}
    the matter ($\Theta_I$), which is invisible because it does not
    source $g_{\mu\nu}$ directly.
  \item \textbf{CTCs}: The imaginary sector \emph{deforms}
    $g_{\mu\nu}$ (via the effective stress-energy $T_{\mu\nu}^{(h)}$)
    to create closed causal curves in the real metric.
\end{itemize}

%----------------------------------------------------------------------
\section{Is an Invisible Object Also Outside the CTC Causal Structure?}
%----------------------------------------------------------------------

In the CTC track, the causal structure is defined by the real metric
$g_{\mu\nu}$ (Definition 2 in \texttt{complex\_time\_causal\_structure.tex}):
\begin{equation}
  \text{causal curve}: \quad g_{\mu\nu}\dot{x}^\mu\dot{x}^\nu \leq 0,
  \quad \operatorname{Re}(\dot\tau) \geq 0.
\end{equation}

An imaginary-sector object $\Theta_I$ with $g_{\mu\nu}[\Theta_I] = 0$:
\begin{itemize}
  \item Has no well-defined worldline in the real metric sense (the
    condition $g_{\mu\nu}\dot{x}^\mu\dot{x}^\nu \leq 0$ is trivially
    satisfied for $g_{\mu\nu} = 0$, but also trivially violated in
    a degenerate way).
  \item Is not constrained by the CTC condition, which requires a
    \emph{non-degenerate} $g_{\mu\nu}$ with compact Killing orbits.
\end{itemize}

\begin{proposition}[Invisibility and CTC conditions are independent]
\label{prop:independent}
In UBT, the invisibility condition ($g_{\mu\nu}[\Theta_I] = 0$)
and the CTC condition ($g_{\mu\nu}$ admits closed causal curves)
are logically independent:
\begin{itemize}
  \item An invisible object ($g_{\mu\nu} = 0$) is neither inside
    nor outside the CTC causal cone, because the causal cone is
    undefined for a degenerate metric.
  \item A CTC-containing spacetime requires $g_{\mu\nu} \neq 0$
    with specific topological properties; invisibility
    ($g_{\mu\nu} = 0$) is incompatible with a CTC-containing
    geometry sourced by the same object.
\end{itemize}
\end{proposition}

%----------------------------------------------------------------------
\section{Coexistence Scenario: CTC Geometry + Invisible Matter}
%----------------------------------------------------------------------

Although an invisible object cannot be a CTC source by itself,
both effects can coexist in a single UBT solution if there are
\emph{two} components: a real-sector source $\Theta_R$ that creates
a CTC-compatible geometry, and an imaginary-sector component $\Theta_I$
that is invisible but gravitates.

In such a solution:
\begin{itemize}
  \item $\Theta_R$ sources $g_{\mu\nu}^{\mathrm{CTC}}$ via the real
    Einstein equations, potentially with CTC-compatible topology.
  \item $\Theta_I$ contributes an additional
    $T_{\mu\nu}[\Theta_I]$ that modifies the CTC geometry through
    back-reaction, possibly expanding or contracting the
    CTC-onset radius (similar to the Kerr–UBT shift derived in
    \texttt{ctc\_conditions.tex}).
  \item Observers in the real sector see: (a) a CTC-compatible
    curvature structure sourced by $\Theta_R$, and (b) an
    additional gravitational field from the invisible $\Theta_I$
    that appears as ``excess curvature'' without a visible source.
\end{itemize}

\begin{conjecture}[Invisible matter modifies CTC geometry]
\label{conj:invisible_ctc}
In a UBT solution with both real and imaginary sector matter,
the imaginary sector shifts the CTC-onset radius by:
\begin{equation}
  r_\mathrm{CTC}^{(\Theta_R + \Theta_I)}
  \;\approx\; r_\mathrm{CTC}^{(\Theta_R)}\left(
    1 + \alpha_I\,\frac{M_I}{M_R}
  \right),
\end{equation}
where $M_I$ and $M_R$ are the effective (gravitational) masses of
the imaginary and real sector components respectively, and $\alpha_I$
is an $O(1)$ dimensionless factor depending on the geometry.
\end{conjecture}

%----------------------------------------------------------------------
\section{Shared Mechanism: Exit from $\operatorname{Re}(\mathcal{G}_{\mu\nu})$}
%----------------------------------------------------------------------

Despite the formal independence of the two phenomena, there is a
conceptually important shared mechanism:

\begin{center}
\begin{tcolorbox}[colback=blue!5, title=Key Insight]
Both \textbf{invisibility} and \textbf{CTCs} in UBT represent forms of
``exit from the observable real metric sector.''
\begin{itemize}
  \item \textbf{Invisibility}: matter exits $\operatorname{Re}(\mathcal{G}_{\mu\nu})$
    by living in $\operatorname{Im}(\mathcal{G}_{\mu\nu})$.
  \item \textbf{CTCs}: causal curves exit the standard causal cone
    of the real metric by exploiting topology in $g_{\mu\nu}$.
\end{itemize}
In both cases, the biquaternionic extension of spacetime provides
the additional structure that makes the phenomenon possible.
Neither effect exists in standard GR (which has only the real metric).
\end{tcolorbox}
\end{center}

This suggests a unifying principle for exotic UBT phenomena:
\emph{the imaginary sector of $\mathcal{G}_{\mu\nu}$ is the locus of
non-standard physics} — dark matter, CTCs, and potentially other
phenomena that are absent in the real (GR) sector.

%----------------------------------------------------------------------
\section{Conclusion}
%----------------------------------------------------------------------

\begin{enumerate}
  \item An imaginary-sector object is \emph{not} inside or outside
    the CTC causal structure — it is \emph{outside the domain of
    definition} of that structure (because the real metric is
    degenerate for an ideal imaginary-sector object).
  \item Invisibility and CTCs are logically \emph{independent}
    phenomena in UBT, using the imaginary sector in complementary ways.
  \item They \emph{can coexist} in a spacetime containing both
    real-sector matter (sourcing CTC geometry) and imaginary-sector
    matter (contributing invisible gravitational fields).
  \item Both phenomena share the conceptual mechanism of ``exit from
    the real metric sector'' — a unifying feature of the UBT
    biquaternionic extension.
  \item The shared geometric origin suggests that dark matter and
    CTC-related phenomenology may be correlated observational signals
    of a non-trivial imaginary metric sector.
\end{enumerate}

\medskip
\noindent\textbf{All statements in this document are speculative.
No claim constitutes an established result of the canonical UBT framework.}

%----------------------------------------------------------------------
\bibliographystyle{plain}
\begin{thebibliography}{9}
\bibitem{CTC_causal}
  D.~Jaroš,
  \emph{ST-1: Causal Structure of Complex Time in UBT},
  Repository document
  \texttt{speculative\_extensions/causality/complex\_time\_causal\_structure.tex}
  (2026).
\bibitem{CTC_conditions}
  D.~Jaroš,
  \emph{ST-3: Closed Timelike Curve Conditions in the Biquaternionic Metric},
  Repository document \texttt{speculative\_extensions/causality/ctc\_conditions.tex}
  (2026).
\bibitem{ST1}
  D.~Jaroš,
  \emph{ST-1: Imaginary Sector Objects in UBT},
  Repository document
  \texttt{speculative\_extensions/invisibility/st1\_imaginary\_object\_definition.tex}
  (2026).
\bibitem{ST2}
  D.~Jaroš,
  \emph{ST-2: Interaction Between Real and Imaginary Sectors in UBT},
  Repository document
  \texttt{speculative\_extensions/invisibility/st2\_sector\_interaction.tex}
  (2026).
\bibitem{ST3}
  D.~Jaroš,
  \emph{ST-3: Stability of Imaginary Sector Configurations in UBT},
  Repository document
  \texttt{speculative\_extensions/invisibility/st3\_stability.tex}
  (2026).
\bibitem{UBT_core}
  D.~Jaroš,
  \emph{Unified Biquaternion Theory: core field equations},
  Repository document \texttt{consolidation\_project/ubt\_core\_main.tex} (2025).
\end{thebibliography}

\end{document}
